% !TEX root = ../main.tex

\chapter{方法设计}

本文面向单视图或三视图鞋履设计草图到三维网格模型的重建任务，设计了一条由草图预处理、草图域适配、三维形状生成和网格质量分析组成的端到端流程。该流程的目标不是直接生成制造级CAD模型，而是在本地硬件约束下生成可预览、可导出、可分析的GLB网格资产，并明确记录生成结果中的几何风险。本章从总体流程、输入表示、ControlNet草图域适配、SF3D基线生成路线、Hunyuan3D生成路线以及网格质量报告六个方面介绍本文的方法设计。

\section{总体流程}

本文系统的核心流程可以概括为：

\begin{enumerate}
  \item 输入单视图或三视图鞋履设计草图；
  \item 对草图进行尺寸归一化、白底合成、灰度化、二值线稿提取和边缘控制图生成；
  \item 根据实验设置选择直接输入路径或ControlNet草图域适配路径；
  \item 将归一化草图或ControlNet生成的中间渲染图送入三维重建后端；
  \item 先以SF3D建立稳定基线和回退路径，再以Hunyuan3D-2mini生成主要形状模型；
  \item 对导出的GLB网格进行非破坏式质量分析，生成面向论文和演示界面的指标报告。
\end{enumerate}

图\ref{fig:method-pipeline}给出了本文方法流程的整体结构。该流程将输入证据、中间图像、三维网格和质量报告都作为可保存结果处理，使论文实验和演示系统能够复用同一套证据链。

\begin{figure}[htbp]
  \centering
  \includegraphics[width=0.96\textwidth]{method/method_pipeline.png}
  \caption{鞋履草图到三维网格模型的总体方法流程}
  \label{fig:method-pipeline}
\end{figure}

在该流程中，预处理模块负责把不同来源、不同尺寸和不同背景的草图整理为统一输入；ControlNet模块负责缓解草图与自然图像或产品渲染图之间的域差异；三维重建模块负责输出可加载的网格模型；网格报告模块负责评价输出文件的基本几何属性和潜在风险。各模块之间通过图像文件、GLB文件和JSON报告连接，使得实验过程中的每个中间结果都可以被保存和复查。

本文设计了两类输入路径。第一类是直接路径，即将归一化后的草图或已有渲染图直接输入三维重建模型。该路径用于检验模型对原始输入的直接理解能力，也适合作为快速推理路径。第二类是ControlNet路径，即先根据草图生成边缘控制图，再通过ControlNet生成更接近产品渲染图的中间图像，最后输入三维重建模型。该路径用于减小设计草图与图像到三维模型训练分布之间的差异。对于实验对比，系统还支持同时运行两条路径，以比较直接草图输入和ControlNet渲染输入对三维结果的影响。

本文的三维重建后端采用“主后端+基线后端”的设计。SF3D保留为稳定基线和回退方案，用于说明快速单图像到三维网格方法在同一输入上的表现；Hunyuan3D-2mini作为当前主要形状生成后端，其优势是在本地环境中可以完成shape-only生成，并且在鞋履样例上的人工视觉审阅结果优于早期SF3D输出。对于三视图输入，本文在流程设计中保留多视图数据组织和融合接口，使侧视图、俯视图和底视图能够作为后续多视图约束或多视图生成模块的输入基础。

\section{输入表示与预处理}

本文正式任务输入为鞋履设计草图。根据设计表达习惯，输入可以是一张侧视或透视草图，也可以是由侧视图、俯视图、底视图等组成的多视图草图。单视图输入操作成本低，适合作为当前本地演示系统的主要路径；三视图输入能够提供更多几何约束，更符合论文题目中“融合三视图特征”的目标，但需要后端模型稳定支持多视图条件。因此，本文在数据组织上保留三视图扩展能力，在当前实现上优先完成单视图可运行闭环。

预处理的目标是将草图整理成三维生成模型和ControlNet模块更容易处理的规范图像。本文采用以下步骤：

\begin{enumerate}
  \item 图像读取与白底合成。输入图像可能带有透明通道或非白色背景，因此首先将图像转为RGBA格式，并与白色背景合成，得到三通道RGB图像。
  \item 方形补边。为了适配扩散模型和三维重建模型常用的方形输入，将原始图像按最长边补成正方形，短边两侧填充白色背景。
  \item 尺寸归一化。将补边后的图像缩放为固定尺寸。当前系统使用$512\times512$作为预处理输出尺寸，以兼顾ControlNet推理效率和图像细节保留。
  \item 灰度化与二值化。将RGB图像转为灰度图，并通过阈值提取线稿区域，得到黑白二值草图。
  \item 边缘控制图生成。使用Canny算子从灰度图中提取边缘，得到ControlNet Canny模式所需的控制图。
  \item Scribble控制图生成。将二值线稿反相，得到亮线条、暗背景的涂鸦式控制图，为后续粗草图输入实验预留接口。
\end{enumerate}

预处理模块输出五类图像：归一化RGB图、灰度图、二值线稿图、Canny控制图和Scribble控制图。其中归一化RGB图用于直接输入路径，Canny或Scribble控制图用于ControlNet路径。保留这些中间结果有两个作用：一是便于定位生成失败来自输入质量、控制图质量还是三维重建模型本身；二是便于在论文和演示界面中展示从草图到三维模型的完整证据链。

对于三视图输入，本文方法要求每个视图独立执行上述预处理，并保留视图名称和视图类型。理想情况下，侧视图主要约束鞋长、鞋面轮廓和鞋底厚度，俯视图主要约束鞋宽、鞋头形状和鞋口范围，底视图主要约束外底形状和底纹区域。当前三视图预处理主要承担输入规范化和视图管理作用，后续可以进一步接入多视图生成或多视图一致性约束模块。

\section{ControlNet草图域适配}

多数开源图像到三维模型更适合处理自然图像或产品渲染图，而不是只有线条和少量标注的设计草图。直接将草图送入三维重建模型时，模型可能把线条理解为平面纹理而非三维结构，导致输出模型过薄、细节缺失或形体塌缩。因此，本文在草图和三维重建后端之间引入ControlNet作为草图域适配模块。

ControlNet通过在预训练扩散模型中加入条件控制分支，使图像生成过程能够遵循边缘、草图、深度等外部控制信号\citep{zhang2023controlnet}。本文使用Canny边缘控制作为默认模式，其输入是预处理阶段生成的边缘图，输出是一张白色背景下的鞋履产品概念渲染图。该中间图像保留草图的鞋底轮廓、鞋面分割、鞋带区域、后跟和鞋头形状，同时补充颜色、明暗和体积感，使其更接近Hunyuan3D或SF3D的常见输入分布。

本文的ControlNet适配过程可以表示为：

\begin{equation}
  I_r = G_{\mathrm{ctrl}}(C(I_s), P, s, \lambda),
\end{equation}

其中，$I_s$表示输入草图，$C(\cdot)$表示边缘或涂鸦控制图生成函数，$P$表示文本提示词，$s$表示随机种子，$\lambda$表示控制强度，$I_r$表示生成的中间渲染图。实际系统中，默认提示词强调“clean side view sneaker product concept render, white background”等约束，以减少背景干扰并突出鞋履主体。

ControlNet模块在本文中不是最终目标模型，而是输入域桥接工具。其优势在于工程实现简单、可解释性较强，并能保存控制图和生成图用于对比；其风险在于生成图可能引入与原始草图不一致的细节，例如改变鞋面纹理、鞋带样式或局部轮廓。因此，本文在实验设计中保留直接路径和ControlNet路径的对比，并将“设计保持程度”作为人工观察的重要内容。

在完成输入域适配后，系统进入三维重建后端选择。图\ref{fig:backend-comparison}展示了SF3D基线路线和Hunyuan3D主路线在本文系统中的功能定位。

\begin{figure}[htbp]
  \centering
  \includegraphics[width=0.96\textwidth]{method/backend_comparison.png}
  \caption{SF3D基线路线与Hunyuan3D主路线的功能定位}
  \label{fig:backend-comparison}
\end{figure}

\section{SF3D基线生成路线}

SF3D是本文保留的稳定基线后端。该方法面向快速三维网格重建，能够从单张图像生成可导出的三维资产\citep{boss2024sf3d}。在本文项目早期实验中，SF3D已经能够从商品鞋图像和鞋履线稿相关输入生成可加载GLB文件，因此适合作为基线模型和系统回退路径。

从方法抽象上看，SF3D所属的前馈式图像到三维网格方法可以表示为从规范化图像到网格资产的直接映射：

\begin{equation}
  \mathcal{M}_{\mathrm{sf}} =
  f_{\theta}^{\mathrm{sf}}(\tilde{I}),
  \quad
  \mathcal{M}_{\mathrm{sf}}=(\mathbf{V}_{\mathrm{sf}}, \mathbf{F}_{\mathrm{sf}}, \mathbf{T}_{\mathrm{sf}}),
  \label{eq:sf3d-forward}
\end{equation}

其中$\tilde{I}$表示预处理后的输入图像，$\mathbf{V}_{\mathrm{sf}}$和$\mathbf{F}_{\mathrm{sf}}$分别表示网格顶点和面片，$\mathbf{T}_{\mathrm{sf}}$表示可选纹理或材质相关信息。与逐样本优化式方法相比，这类前馈映射的优势是推理流程短、工程封装简单，适合为本文系统建立可复现基线。

如果从图像到网格学习的一般原理理解，前馈式三维重建通常需要同时约束投影视觉一致性和几何合理性，可抽象为如下组合目标：

\begin{equation}
  \mathcal{L}_{\mathrm{mesh}} =
  \lambda_s \mathcal{D}_{\mathrm{sil}}\left(\Pi(\mathcal{M}), I\right)
  + \lambda_n \mathcal{D}_{\mathrm{normal}}\left(N(\mathcal{M}), N^\ast\right)
  + \lambda_g \mathcal{R}_{\mathrm{geo}}(\mathcal{M}),
  \label{eq:mesh-objective}
\end{equation}

其中$\Pi(\cdot)$表示渲染或投影过程，$\mathcal{D}_{\mathrm{sil}}$表示轮廓或外观差异，$\mathcal{D}_{\mathrm{normal}}$表示法线一致性约束，$\mathcal{R}_{\mathrm{geo}}$表示平滑、连通性或形状先验等几何正则项。本文引用该公式并非复现SF3D训练过程，而是说明快速图像到网格方法为什么能够在有限输入下输出三维形状：模型在训练阶段已经学习了图像观测与三维几何之间的统计对应关系，推理阶段则将这种先验压缩为一次前向计算。

SF3D路线与Hunyuan3D路线共享同一套输入预处理和ControlNet适配结果。对于同一输入，系统可以选择将归一化草图或ControlNet渲染图送入SF3D，并导出GLB模型。由于SF3D运行稳定、输出文件结构清晰，它可以用于验证整个流水线的基本可用性，也可以在Hunyuan3D后端不可用时保持系统仍能输出三维模型。

本文没有将SF3D作为主后端，主要原因在于鞋履样例的人工审阅结果显示其内部结构和局部几何合理性不足。尤其是对于线稿或单视图输入，SF3D可能生成较薄的几何体或包含较多碎片，难以表现鞋履应有的厚度和内部空间。因此，本文将SF3D定位为“稳定基线”而非“最佳结果模型”。在实验章节中，SF3D结果将用于说明ControlNet渲染输入相对于直接草图输入的改善，以及Hunyuan3D-2mini相对于早期基线的优势。

\section{Hunyuan3D生成路线}

Hunyuan3D-2系列是本文当前主要三维生成路线的基础。该系列模型面向图像到三维资产生成，包含形状生成和纹理生成等模块\citep{hunyuan3d2025hunyuan3d2}。考虑到论文任务更关注鞋履形状而非纹理质量，本文当前采用Hunyuan3D-2mini的shape-only路径作为主要生成方式。

Hunyuan3D路线的输入可以是归一化草图，也可以是ControlNet生成的鞋履渲染图。系统调用Hunyuan3D推理脚本后，模型生成三维形状并导出GLB网格文件。为了降低显存压力，当前实现采用fp16模型权重、固定推理步数、指定octree resolution和分块参数，并关闭纹理生成。经过本地验证，Hunyuan3D-2mini能够在约4.4GB峰值CUDA显存下完成鞋履样例的shape-only生成，满足本文本地演示系统的基本需求。

从条件生成角度看，Hunyuan3D路线可以抽象为图像条件下的三维形状生成过程。给定输入图像特征$\phi(\tilde{I})$，模型在形状潜空间中从噪声变量逐步恢复三维表示：

\begin{equation}
  z_T \sim \mathcal{N}(0, I), \quad
  z_{t-1} = D_{\theta}(z_t, t, \phi(\tilde{I})),
  \quad
  \mathcal{M}_{\mathrm{hy}} = \mathcal{E}(z_0),
  \label{eq:hunyuan-generation}
\end{equation}

其中$z_t$表示第$t$步的形状潜变量，$D_{\theta}$表示条件去噪或生成网络，$\mathcal{E}$表示从最终形状表示到网格的提取或导出过程。该表达强调了Hunyuan3D与前馈基线的区别：它不是只依赖一次图像到网格映射，而是利用三维生成先验在条件约束下补全不可见形状，因此更适合承担本文的主要外观级三维生成任务。

本文选择Hunyuan3D-2mini作为主后端有三个原因。首先，它在本地硬件上可运行，不需要高显存服务器即可完成主要演示流程。其次，在项目人工审阅中，它生成的鞋履外观明显优于早期SF3D结果，更适合作为当前论文和演示的主路线。再次，它属于较新的开源三维生成模型，便于在论文中与扩散先验和图像到三维生成研究联系起来。

但本文也明确限定Hunyuan3D路线的适用边界。第一，当前系统使用的是shape-only生成，不宣称已经解决高质量纹理重建。第二，Hunyuan3D-2mini主要使用单图像输入，不等同于真正完成了三视图特征融合。第三，模型生成的三维形状虽然外观更自然，但仍然无法从单视图可靠推断鞋履内部结构，例如鞋口内部空间、鞋帮厚度和真实穿着腔体。因此，本文将Hunyuan3D路线定位为当前最佳本地外观级形状生成方案，而不是完整鞋履结构建模方案。

\section{网格后处理与质量报告}

三维生成模型输出的GLB文件需要经过质量检查才能用于论文分析和演示展示。本文采用非破坏式网格报告作为当前后处理策略，即只读取和分析网格，不默认执行几何修改。这样可以避免自动平滑、减面或修复操作破坏鞋履设计细节，也能保证实验记录反映模型原始输出质量。

网格报告模块使用trimesh加载GLB文件，并提取以下指标：

\begin{itemize}
  \item 文件存在性和文件大小，用于判断导出是否成功；
  \item 场景类型和几何体数量，用于判断GLB内部结构是否可解析；
  \item 顶点数和面数，用于衡量网格复杂度；
  \item 包围盒最小点、最大点和三维尺寸，用于描述模型空间尺度；
  \item 最小维度与最长维度之比，作为厚度或扁平化风险的代理指标；
  \item 连通组件数量和最大组件面数，用于发现碎片化几何；
  \item 表面积、水密性和可靠体积，用于辅助判断网格封闭性；
  \item 结构化警告，用于提示薄片化、碎片化、非水密和内部结构不可验证等风险。
\end{itemize}

其中，厚度代理指标定义为：

\begin{equation}
  r_t = \frac{\min(d_x, d_y, d_z)}{\max(d_x, d_y, d_z)},
\end{equation}

其中$d_x$、$d_y$、$d_z$表示模型包围盒在三个方向上的尺寸。若$r_t$过小，说明模型在某一方向上明显偏薄，可能存在扁平化或深度不足问题。该指标不能直接等价于真实鞋宽或鞋底厚度，但可以作为筛查异常几何的简明信号。

内部结构风险是本文特别保留的警告项。即使一个GLB模型水密、连通组件数量较少、外观轮廓合理，自动网格指标仍无法证明鞋履内部结构符合真实鞋类设计要求。因为内部鞋腔、鞋帮厚度和穿着空间往往不可从单视图观察得到，也很难仅通过包围盒和面数判断。因此，本文的网格报告总是保留\texttt{internal\_structure\_not\_validated}警告，提醒使用者需要通过人工剖视、更多视图或结构先验进一步验证。

综上，本文方法采用“生成模型输出+非破坏式质量报告”的组合方式。该设计既能展示当前开源三维生成模型在鞋履草图任务上的可用性，又能避免对结果质量进行过度承诺。若引入几何修复、平滑或减面模块，也应在保留原始GLB的基础上进行对比，并经过人工视觉审阅后再作为默认输出。
